\documentclass[10pt]{article}
\usepackage[margin=12mm]{geometry}
\usepackage{graphicx}
\usepackage{caption}
\usepackage{tikz}
\pagestyle{empty}
\captionsetup{font=small,labelfont=bf}

\begin{document}
\begin{figure}[p]
\centering
\resizebox{\textwidth}{!}{\input{fig3A_ce_null_models.tex}}

\vspace{4mm}
\begin{minipage}[t]{0.485\textwidth}
  \centering
  \includegraphics[width=\linewidth]{fig3B_ce_edge_retention_tree_distance.pdf}
\end{minipage}\hfill
\begin{minipage}[t]{0.485\textwidth}
  \centering
  \includegraphics[width=\linewidth]{fig3C_ce_structural_retention.pdf}
\end{minipage}

\caption{\textbf{Null models and structural metrics along the terminal-cell Pareto front.}
\textbf{(A)} Nested cousin-shuffling null models progressively relax the
lineage restriction on assignments of terminal cells to candidate parents. First-cousin
shuffling acts within the smallest neighborhoods, second- and third-cousin
shuffling act within successively larger ancestral neighborhoods, and fully random
shuffling permits assignments among all candidate parents. Terminal-cell
identities, defined by their final physical positions and molecular states,
remain fixed to strictly follow the biological outcome of embryogenesis.
\textbf{(B)} Edge retention (blue, left axis) and mean lineage-tree distance
among changed edges (grey, inverted right axis) compared with the natural
lineage for all assignments along the Pareto front. The natural lineage has
edge retention one and, by definition, no changed edges; it is shown at the
zero-distance reference. Markers identify the natural lineage, the
maximum-edge-retention assignment, and the minimum-tree-distance assignment.
Marks on the right axis show the corresponding mean tree distances among changed edges
under first-cousin and full-random shuffling. Approaching either single-objective
optimum requires increasingly large changes from the natural lineage.
\textbf{(C)} Structural retention while moving along the Pareto front from
the maximum-edge-retention compromise toward either single-objective
optimum. The inset indicates the direction of travel of the two colored
trajectories along the Pareto front. For each trajectory separately, the
horizontal axis reports the fraction of the theoretically attainable reduction in its
corresponding distance objective; this normalization does not compare or combine the
objectives. Natural-edge retention falls sharply as the remaining benefit is pursued.
The annotations report changes over the shaded final 5\% of attainable
distance reduction, using mean lineage-tree distance across all terminal
edges: 5.4 percentage points of edge retention and 0.7 tree-distance units
toward the travel optimum, and 11.4 percentage points and 2.0 tree-distance
units toward the cell-state optimum.}
\label{fig:ce_terminal_pareto_supporting}
\end{figure}
\end{document}
